% sections/03-generacion-imagen.tex — tarea «Generación Imagen Forense».
% El enunciado detallado de esta parte no venía en la guía transcrita
% (docs/ws-04.md llega hasta el punto 5.2); la estructura queda montada y las
% cajas vacías delatan lo que falta cuando se compile con FINAL=1.
\section{Generación de la imagen forense}

La adquisición de una imagen forense consiste en producir una copia bit a bit del medio de almacenamiento (no una copia de archivos, que perdería el espacio no asignado y con él los artefactos borrados), acompañada de una verificación por función \eng{hash} que permita demostrar ante un tribunal que lo analizado es idéntico a lo incautado \cite{nist80086}; el procedimiento se ejecuta sobre el medio original en modo de solo lectura, preferiblemente a través de un bloqueador de escritura, y tanto el \eng{hash} del origen como el de la imagen resultante se consignan en la cadena de custodia.

\subsection{Identificación del medio y adquisición}

\begin{lstlisting}[language=bash]
sudo fdisk -l                      # identificar el dispositivo de origen
sudo dd if=/dev/sdX of=evidencia.dd bs=4M status=progress conv=noerror,sync
sudo dcfldd if=/dev/sdX of=evidencia.dd hash=sha256 hashlog=evidencia.hash
\end{lstlisting}

\captura{imagen-identificacion-medio.png}{Identificación del medio de almacenamiento que se va a adquirir.}

\begin{analisis}
\end{analisis}

\captura{imagen-adquisicion.png}{Proceso de adquisición de la imagen forense en curso.}

\begin{analisis}
\end{analisis}

\subsection{Verificación de integridad}

\begin{lstlisting}[language=bash]
sha256sum /dev/sdX
sha256sum evidencia.dd
\end{lstlisting}

\captura{imagen-hash-verificacion.png}{Comparación de los \eng{hashes} del medio original y de la imagen obtenida.}

\begin{analisis}
\end{analisis}

\pregunta{¿Por qué la imagen forense debe acompañarse de su función \eng{hash} y qué ocurre si los valores del origen y de la copia no coinciden?}
\begin{respuesta}
\end{respuesta}
